\chapter{GT-TİPİ YENİDEN YAPILANDIRILABİLİR ENERJİ DEPOLAMA MİMARİSİ}
\label{chap:gttipi}

Bölüm~\ref{chap:anahtarlama}--\ref{chap:deney}, faz kolu ölçeğinde tezin ilk iki katkısını kurmuştur: T-tipi dönüştürücünün düşük güç davranışı, ters akım dolaşımı, nötr anahtar susturma yöntemi, Sin\textsuperscript{3} tabanlı ileri besleme kompanzasyonu ve deney düzeneği incelenmiştir. Bölüm~\ref{chap:gttipi} sistem ölçeğine geçerek tezin üçüncü katkısını, aynı T-tipi yarı iletken tabanının enerji depolama sisteminin şarj ve deşarj işlevlerinde yeniden görevlendirilmesini ele almaktadır. Anma görünür gücü \SI{20}{\kilo\volt\ampere}, tasarım güç aralığı \num{20}--\SI{25}{\kilo\watt} olan deney düzeneği GT-tipi olarak kurulmuştur: Şekil~\ref{fig:gelen-tam-topoloji} üzerindeki üç fazlı güç topolojisinde her faz hücresindeki röle ve doğrultucu diyotlar kip değişimini, alt bölümdeki transformatörler ise şarj yönündeki batarya yönetim sistemi (BYS) transformatörlerini oluşturmaktadır.

``GT-tipi'' adı, tezde \emph{geliştirilmiş/değiştirilmiş T-tipi} enerji depolama mimarisini belirtmek için kullanılan bir adlandırmadır; literatürde standartlaşmış bağımsız bir çok seviyeli dönüştürücü ailesi değildir. Mimari; T-tipi faz kollarını, kip seçen RL1 fonksiyon anahtarını, nötr IGBT'lerle kurulan Vienna PFC katını, çok sekonderli BYS transformatörlerini ve batarya başına doğrultucu--buck şarj kollarını tek bir donanımda birleştirmektedir.

Bölüm boyunca deney sonucu, benzetim kaydı ve tasarım tanımı birbirinden ayrılmıştır. Evirici kipinin düşük güç davranışı Bölüm~\ref{chap:deney} altındaki deneylerle, şarj yönünün işlevsel bütünlüğü Alt Bölüm~\ref{sec:gt-benzetim} altındaki benzetim kayıtlarıyla gösterilmiştir. Çalışma kipleri, geçiş kuralları ve koruma işlevleri ise tasarım tanımı olarak verilmiştir. Devre şemaları ve benzetim ekranları tez sahibinin tasarım çalışmalarından alınmış ve özgün çözünürlükleriyle verilmiştir.

\section{Tasarım Gereksinimleri}
\label{sec:gt-gereksinimler}

Yüksek güçlü bir EDS'de batarya dizisi sabit ve ideal bir DA kaynak değildir. Hücre veya modül uç gerilimleri şarj durumu, sıcaklık, iç direnç, akım ve yaşlanmayla değişir; seri dizinin kullanılabilir kapasitesi zayıf birim tarafından sınırlandırılabilir \citep{Plett2015,Hannan2017,Farakhor2023}. Yeniden yapılandırılabilir batarya sistemleri, seçilen modülleri enerji yoluna alma veya yoldan çıkarma olanağıyla arıza yalıtımı, gerilim uyarlama ve kapasite kullanımını iyileştirmeyi hedefler \citep{Momayyezan2016,Farakhor2023,Huang2024}. İlave yarı iletkenler ise iletim kaybı, sürücü sayısı, koruma gereksinimi ve denetim karmaşıklığı oluşturduğundan tasarım değerlendirmesi yalnız işlev sayısına değil, kayıp ve güvenlik bedeline de dayanmalıdır.

GT-tipi mimarinin güç tabanı Bölüm~\ref{chap:sistem} altında tanımlanan deney düzeneğidir: üç fazlı T-tipi evirici, bölünmüş DA bara, LCL filtre ve \SI{20}{\kilo\volt\ampere} anma görünür gücü (\num{20}--\SI{25}{\kilo\watt} tasarım güç aralığı). Mimarinin amacı ayrı bir şarj dönüştürücüsü eklemek yerine aynı faz kollarını şarj yönünde de kullanmaktır; gerilim, akım ve anahtarlama frekansı sınırları evirici kipiyle ortaktır ve Tablo~\ref{tab:sistem-parametreleri} içinde verilmiştir.

Tasarım gereksinimleri beş işlevsel ilkeye indirgenmiştir: enerji yolunun çift yönlü olabilmesi, şarj sırasında şebeke akım kalitesinin gözetilmesi, batarya kollarının galvanik olarak ayrılması, her bataryanın güvenli akım--gerilim--sıcaklık sınırlarında tutulması ve her kip geçişinin ölçüm geri bildirimiyle güvenli biçimde tamamlanması. IEC~62619 endüstriyel ve sabit uygulamalardaki ikincil lityum hücre ve bataryaların güvenli işletimine; IEC~62477-1 güç elektroniği dönüştürücülerindeki elektrik çarpması, yangın, ısıl ve enerji tehlikelerine; IEC~62933-5-2 ise şebeke ile bütünleşik elektrokimyasal EDS'nin yaşam döngüsü boyunca alt sistem etkileşimlerinden doğan güvenlik risklerine odaklanmaktadır \citep{IEC62619,IEC62477,IEC62933_5_2}.

\section{GT-Tipi Topolojinin Yapısı}
\label{sec:gt-yapi}

GT-tipi yapı, deney düzeneğindeki T-tipi faz kollarını RL1 fonksiyon anahtarı ve doğrultucu/boost diyotlarıyla üç işlev arasında yeniden görevlendirmektedir: şebekeye güç veren evirici, BYS transformatörlerini besleyen iki seviyeli yarım köprü ve şebekeden şarj akımı çeken Vienna PFC. T-tipi faz kolunun temel yapısı Bölüm~\ref{chap:sistem} içindeki Şekil~\ref{fig:t-tipi-uc-seviye} ve Şekil~\ref{fig:t-tipi-akim-yollari} ile verildiğinden yeniden basılmamıştır. GT-tipi yapının ayırt edici alt devreleri Şekil~\ref{fig:gt-viyana-pfc}--\ref{fig:gt-pfc-modu} üzerinde sırasıyla açıklanmaktadır. Elemanların iki enerji yönündeki görevleri Tablo~\ref{tab:gt-eleman-gorevleri} üzerinde eşlenmiştir.

\begin{table}[H]
  \centering
  \caption{GT-tipi yapıda faz kolu elemanlarının deşarj ve şarj yönündeki görevleri.}
  \label{tab:gt-eleman-gorevleri}
  \small
  \begin{tabularx}{\textwidth}{@{}>{\raggedright\arraybackslash}p{0.17\textwidth}>{\raggedright\arraybackslash}X>{\raggedright\arraybackslash}X@{}}
    \toprule
    Eleman & Deşarj (evirici kipi) & Şarj (PFC ve BYS kipi) \\
    \midrule
    $Q_1$--$Q_2$ & T-tipi kolun dış anahtarları, \SI{10}{\kilo\hertz} PWM & İki seviyeli yarım köprü, BYS transformatörü primer beslemesi \\
    $Q_3$--$Q_4$ & Nötr kol, $O$ durumu ve düşük güçte nötr anahtar susturma & Vienna PFC kolunun nötr anahtarları \\
    $D_1$--$D_2$ & Kararlı çalışmada iletimde değil & PFC yükseltici diyotları \\
    RL1 & Evirici konumu: T-tipi kol ve LCL filtre & İki seviyeli yüksek frekans katı ve ayrı PFC yolu konumu \\
    $L_1$--$C_f$--$L_2$ & LCL çıkış filtresi & Vienna PFC giriş süzgeci \\
    \bottomrule
  \end{tabularx}
\end{table}

Şekil~\ref{fig:gt-viyana-pfc} üzerindeki Vienna doğrultucu; üç faz giriş endüktansları, altı doğrultma diyodu, nötr IGBT grubu ve bölünmüş DA baradan oluşmaktadır. Şarj yönünde şebekeden düşük harmonikli ve birim güç faktörüne yakın akım çekmek için Vienna katı kullanılmaktadır.

\begin{figure}[H]
  \centering
  \includegraphics[width=0.92\textwidth]{gt-tipi/onceki_tez_6_2_viyana_pfc.jpeg}
  \caption{Vienna doğrultucu PFC devresi.}
  \label{fig:gt-viyana-pfc}
\end{figure}

Şekil~\ref{fig:gt-cok-sekonderli-bys}, tek primer ve çok sayıda sekonder sargı üzerinden her bataryaya ayrı doğrultucu--buck kolu sağlamaktadır. Her çıkışta diyot doğrultma, DA bağlantı kondansatörü, kontrollü anahtar, serbest dolaşım diyodu, endüktans ve çıkış kondansatörü bulunmaktadır. Batarya bazında şarj akımı ve gerilimi ayarlama hedefi sekonder kollarının fiziksel ayrışmasına dayanmaktadır.

\begin{figure}[H]
  \centering
  \includegraphics[width=0.67\textwidth]{gt-tipi/onceki_tez_6_3_cok_sekonderli_bys.jpeg}
  \caption{Tek primer ve çoklu sekonder sargılı transformatör içeren aktif BYS prensip devresi.}
  \label{fig:gt-cok-sekonderli-bys}
\end{figure}

Şekil~\ref{fig:gt-ana-topoloji}, $Q_1$--$Q_2$ ile $Q_3$--$Q_4$ anahtarlarının ortak T-tipi çıkış düğümündeki bağlantısını RL1 konumuna göre değiştiren ana topolojiyi göstermektedir (Tablo~\ref{tab:gt-eleman-gorevleri}).

\begin{figure}[H]
  \centering
  \includegraphics[width=0.82\textwidth]{umit_2026-11-09/gt_tipi_ana_topoloji_rl1.jpg}
  \caption{GT-tipi ana topoloji ve RL1 fonksiyon anahtarı.}
  \label{fig:gt-ana-topoloji}
\end{figure}

Şekil~\ref{fig:gt-evirici-lcl}, RL1 ve RL2'nin evirici konumunda $Q_1$--$Q_4$ grubunun üç seviyeli T-tipi kolu, $L_1$--$C_f$--$L_2$ ağının ise şebeke akımı filtresini oluşturduğu deşarj yolunu göstermektedir.

\begin{figure}[H]
  \centering
  \includegraphics[width=0.76\textwidth]{umit_2026-11-09/gt_tipi_evirici_lcl_rl2.jpg}
  \caption{GT-tipi topolojinin evirici çalışma modu ve LCL filtre bağlantısı.}
  \label{fig:gt-evirici-lcl}
\end{figure}

Şekil~\ref{fig:gt-iki-seviye-trafo}, şarj yönünde üç fazın $Q_1$--$Q_2$ yarım köprülerinin üç çok sekonderli BYS transformatörünün primerlerini beslediği bağlantıyı göstermektedir.

\begin{figure}[H]
  \centering
  \includegraphics[width=0.88\textwidth]{gt-tipi/onceki_tez_6_6_iki_seviye_trafo_modu.jpeg}
  \caption{İki seviyeli evirici çıkışları ve çok sekonderli BYS transformatörleri.}
  \label{fig:gt-iki-seviye-trafo}
\end{figure}

Şekil~\ref{fig:gt-pfc-modu}, nötr IGBT çiftlerinin PFC yükseltici diyotları ($D_1$--$D_2$) ve LCL elemanlarıyla üç fazlı Vienna PFC kollarına dönüştürüldüğü şarj yolunu göstermektedir.

\begin{figure}[H]
  \centering
  \includegraphics[width=0.94\textwidth]{gt-tipi/pfc_modu_q1q3_q2q4.png}
  \caption{Nötr IGBT gruplarıyla oluşturulan üç fazlı PFC bölümü. Her fazda $Q_4$ DA bara orta noktası (NÖTR) tarafındaki, $Q_3$ faz düğümü tarafındaki nötr IGBT'dir.}
  \label{fig:gt-pfc-modu}
\end{figure}

Şarjda her batarya ayrı izole buck kolundan bağımsız akım ve gerilim düzenlemesiyle beslenmekte, deşarjda ise enerji ortak DA bara ve T-tipi evirici üzerinden şebekeye aktarılmaktadır. Batarya bazında dengeleme hedefi, şarj yönündeki batarya başına düzenlemeyle sağlanmaktadır.

\begin{figure}[H]
  \centering
  \begin{tikzpicture}[
    font=\small,
    block/.style={draw,rounded corners,align=center,minimum width=26mm,minimum height=12mm,fill=blue!5},
    batt/.style={draw,rounded corners,align=center,minimum width=26mm,minimum height=12mm,fill=green!8},
    ctrl/.style={draw,rounded corners,align=center,minimum width=29mm,minimum height=12mm,fill=orange!10},
    flow/.style={-{Latex[length=2.2mm]},thick},
    signal/.style={-{Latex[length=2mm]},dashed}
  ]
    \node[block] (grid) {Şebeke\\LCL filtre};
    \node[block,right=9mm of grid] (pfc) {Nötr IGBT'li\\Vienna PFC};
    \node[block,right=9mm of pfc] (dc) {Bölünmüş\\DA bara};
    \node[block,right=9mm of dc] (hb) {$Q_1$--$Q_2$\\yarım köprü};
    \node[block,below=12mm of hb] (tr) {Çok sekonderli\\BYS trafoları};
    \node[batt,left=9mm of tr] (buck) {Batarya başına\\doğrultucu--buck};
    \node[batt,left=9mm of buck] (battery) {Batarya\\grupları};
    \node[block,above=12mm of pfc] (inv) {$Q_1$--$Q_4$\\T-tipi evirici};
    \node[ctrl,below=12mm of buck] (ems) {RL1 kip seçimi ve\\üst seviye denetim};

    \draw[flow] (grid) -- node[above,font=\scriptsize]{şarj} (pfc);
    \draw[flow] (pfc) -- (dc);
    \draw[flow] (dc) -- (hb);
    \draw[flow] (hb) -- (tr);
    \draw[flow] (tr) -- (buck);
    \draw[flow] (buck) -- (battery);
    \draw[flow] (dc) |- node[pos=0.25,right,font=\scriptsize]{deşarj} (inv);
    \draw[flow] (inv) -| (grid);
    \draw[signal] (ems) -| (pfc);
    \draw[signal] (ems.east) -| ([xshift=7mm]tr.east) |- (hb.east);
    \draw[signal] (ems) -- (buck);
  \end{tikzpicture}
  \caption{GT-tipi enerji ve denetim katmanları; düz oklar şebekeden şarj ve şebekeye deşarj yönlerini, kesikli oklar kip seçimi ile denetim işaretlerini göstermektedir.}
  \label{fig:gt-sentez-mimarisi}
\end{figure}

Şekil~\ref{fig:gt-sentez-mimarisi}, farklı şekillere dağılan güç ve denetim katmanlarını tek bir işlevsel akışta birleştirmektedir. Şarj yönünde PFC katı şebeke akımı ile DA bara gerilimini, yarım köprü ve transformatör katı yalıtılmış enerji aktarımını, batarya başına buck kolları şarj akım ve gerilimlerini belirler. Deşarj yönünde T-tipi evirici, Bölüm~\ref{chap:sin3} altındaki akım denetimiyle şebekeye güç verir. Aynı anda bütün katların hızlı anahtarlanması gerekmez; RL1 kip geçişleri, güç dönüştürücü PWM'ninden daha yavaş bir üst seviye durum makinesiyle yönetilmelidir.

\section{Şarj Yönündeki Benzetim Kayıtları}
\label{sec:gt-benzetim}

GT-tipi topolojinin şarj yönündeki işlevsel bütünlüğü; Vienna PFC, iki seviyeli transformatör besleme katı ve batarya şarj kolunu aynı Simulink modeli içinde çalıştıran bir benzetimle gösterilmiştir. Benzetim, Bölüm~\ref{chap:anahtarlama}--\ref{chap:deney} içindeki evirici kipine ait S1--S4 zincirinden ayrı tutulmuş ve mimari tanımın hemen ardından verilmiştir. Şekil~\ref{fig:gt-simulink-modeli}--\ref{fig:gt-batarya-soc}, şebekeden bataryaya uzanan PFC $\rightarrow$ DA bara $\rightarrow$ transformatör primeri $\rightarrow$ sekonder $\rightarrow$ buck/batarya zincirinin her aşamasını dalga biçimleriyle belgelemektedir.

Şekil~\ref{fig:gt-simulink-modeli}, PFC denetimi, DA bara, üç fazlı dönüştürücü, çok sekonderli transformatörler ve batarya alt sistemlerinin aynı modelde kurulduğunu göstermektedir.

\begin{figure}[H]
  \centering
  \includegraphics[width=0.98\textwidth]{gt-tipi/onceki_tez_6_8_simulink_modeli.png}
  \caption{Şarj yönündeki Vienna PFC, iki seviyeli evirici ve transformatör tabanlı BYS Simulink modeli.}
  \label{fig:gt-simulink-modeli}
\end{figure}

Şekil~\ref{fig:gt-pfc-akimlari}, PFC sonrasında üç faz akımının dengeli ve yaklaşık sinüzoidal biçimde oluştuğunu göstermektedir.

\begin{figure}[H]
  \centering
  \includegraphics[width=0.96\textwidth]{gt-tipi/onceki_tez_6_9_pfc_sebeke_akimlari.jpeg}
  \caption{Vienna PFC sonrasındaki üç faz şebeke akımları.}
  \label{fig:gt-pfc-akimlari}
\end{figure}

Şekil~\ref{fig:gt-da-bara}, DA bara gerilimlerinin başlangıç geçicisinden sonra birbirine yakın bir çalışma bandına yerleştiğini göstermektedir.

\begin{figure}[H]
  \centering
  \includegraphics[width=0.96\textwidth]{gt-tipi/onceki_tez_6_10_da_bara_gerilimleri.png}
  \caption{DA bara gerilimlerinin zamanla değişimi.}
  \label{fig:gt-da-bara}
\end{figure}

İki seviyeli yüksek frekans katının transformatör primerine uyguladığı gerilim Şekil~\ref{fig:gt-primer-gerilim}, primer akımı Şekil~\ref{fig:gt-primer-akim} üzerindedir. Primer geriliminin pozitif, sıfır ve negatif seviyeler arasında değişmesi, görev oranı üzerinden transformatöre aktarılan gücün ayarlanabildiğini gösterir. Primer akımı aynı anahtarlama dizisine yanıt vermekte; kenarlardaki eğimler eşdeğer endüktans ve manyetik devre dinamiğinin etkisini göstermektedir.

\begin{figure}[H]
  \centering
  \begin{subfigure}[t]{0.48\textwidth}
    \centering
    \includegraphics[width=\linewidth]{gt-tipi/onceki_tez_6_11_primer_gerilim.png}
    \caption{Primer sargı gerilimi}
    \label{fig:gt-primer-gerilim}
  \end{subfigure}\hfill
  \begin{subfigure}[t]{0.48\textwidth}
    \centering
    \includegraphics[width=\linewidth]{gt-tipi/onceki_tez_6_12_primer_akim.png}
    \caption{Primer akımı}
    \label{fig:gt-primer-akim}
  \end{subfigure}
  \caption{BYS transformatörü primer sargı gerilimi ve akımı.}
  \label{fig:gt-primer}
\end{figure}

Şekil~\ref{fig:gt-sekonder}, bir bataryayı besleyen sekonder sargının doğrultma öncesi gerilim ve akımını göstermektedir. Gerilim işaret değiştiren darbeli yapıya sahipken akımın pozitif enerji aktarım aralıklarında belirginleşmesi, sekonder doğrultucu ve buck girişinin primer anahtarlamasıyla eşzamanlı çalıştığını gösterir.

\begin{figure}[H]
  \centering
  \begin{subfigure}[t]{0.48\textwidth}
    \centering
    \includegraphics[width=\linewidth]{gt-tipi/onceki_tez_6_13_sekonder_gerilim.png}
    \caption{Sekonder sargı gerilimi}
    \label{fig:gt-sekonder-gerilim}
  \end{subfigure}\hfill
  \begin{subfigure}[t]{0.48\textwidth}
    \centering
    \includegraphics[width=\linewidth]{gt-tipi/onceki_tez_6_14_sekonder_akim.png}
    \caption{Sekonder sargı akımı}
    \label{fig:gt-sekonder-akim}
  \end{subfigure}
  \caption{BYS transformatörü sekonder sargı gerilimi ve akımı.}
  \label{fig:gt-sekonder}
\end{figure}

Enerji zincirinin batarya tarafındaki son kaydı Şekil~\ref{fig:gt-batarya-soc} üzerindedir. Kayıtta bir bataryanın şarj durumu (State of Charge, SoC), uç gerilimi ve akımı aynı zaman ekseninde verilmiştir.

\begin{figure}[H]
  \centering
  \includegraphics[width=0.80\textwidth]{gt-tipi/onceki_tez_6_15_batarya_soc.jpeg}
  \caption{Bir bataryaya ait SoC, uç gerilimi ve akım kayıtları.}
  \label{fig:gt-batarya-soc}
\end{figure}

Şekil~\ref{fig:gt-simulink-modeli}--\ref{fig:gt-batarya-soc} kayıtları, şarj yönünün işlevsel bütünlüğünü nitel olarak göstermektedir. Alt Bölüm~\ref{sec:gt-dogrulama} altındaki doğrulama matrisi, iki enerji yönünün kanıtlarını katman katman özetlemektedir.

\section{Şarj ve Deşarj Çalışma Modları}
\label{sec:gt-modlar}

GT-tipi yapının çalışması tek bir ``çift yönlü'' komutla açıklanamaz. RL1 konumu, PFC ya da evirici olarak çalışan faz kolları, BYS transformatörlerini besleyen yarım köprüler, batarya buck kolları ve kontaktörler için izin verilen durumların birlikte tanımlanması gerekir. Tablo~\ref{tab:gt-calisma-modlari}, Şekil~\ref{fig:gt-ana-topoloji}--\ref{fig:gt-pfc-modu} üzerindeki işlevleri güvenli bir durum dizisine dönüştürmektedir.

\begin{table}[H]
  \centering
  \caption{GT-tipi mimarinin önerilen çalışma modları ve geçiş koşulları.}
  \label{tab:gt-calisma-modlari}
  \small
  \begin{tabularx}{\textwidth}{@{}p{0.18\textwidth}p{0.27\textwidth}X@{}}
    \toprule
    Mod & Güç yolu & Temel koşul ve amaç \\
    \midrule
    Güvenli kapalı & Bütün PWM'ler ve ana kontaktörler kapalı & Enerjisiz başlatma, kilitli hata veya bakım durumu; artık gerilim denetlenir. \\
    Ön şarj / öz sınama & Sınırlı akımlı DA bara yolu ve ölçümler & İzolasyon, kontaktör ve röle geri bildirimi, sensör makullüğü ve DA bara gerilim yükselme hızı kontrol edilir. \\
    Şebekeden şarj & Şebeke--LCL--Vienna PFC--DA bara--$Q_1$--$Q_2$ yarım köprü--BYS trafosu--buck--batarya & Şebeke akımı, DA bara, transformatör ve batarya sınırları birlikte korunur. \\
    CC şarj & Batarya başına buck akım düzenleme & Batarya üreticisinin akım ve sıcaklık sınırları içinde sabit akım uygulanır. \\
    CV şarj & Batarya başına buck gerilim düzenleme & Üst gerilim sınırında akım azalıncaya kadar sabit gerilim uygulanır. \\
    Şebekeye deşarj (evirici kipi) & DA bara--$Q_1$--$Q_4$ T-tipi kol--LCL--şebeke & Şebeke akımı denetlenir; düşük güçte nötr anahtar susturma ve Sin\textsuperscript{3} ileri besleme kullanılabilir. \\
    RL1 kip geçişi & PWM durdurulmuş, röle konum değiştiriyor & Faz akımları sıfıra yakın ve DA bara yarıları dengedeyken röle konumu değiştirilir; geri bildirim doğrulanmadan yeni kip başlatılmaz. \\
    \bottomrule
  \end{tabularx}
\end{table}

Batarya başına şarj referansı, CC--CV geçişini açıkça göstermek için
\begin{equation}
  i_{b,k}^{\star}=
  \begin{cases}
    I_{\mathrm{CC},k}, & v_{b,k}<V_{\mathrm{CV},k}\ \land\ T_{\min,k}<T_{b,k}<T_{\max,k},\\
    \operatorname{sat}\!\left[K_{v,k}\left(V_{\mathrm{CV},k}-v_{b,k}\right),0,I_{\mathrm{CC},k}\right],
      & v_{b,k}\geq V_{\mathrm{CV},k},\\
    0, & \text{koruma koşulu}
  \end{cases}
  \label{eq:gt-cccv-referansi}
\end{equation}
biçiminde tanımlanabilir. Eşitlik~\eqref{eq:gt-cccv-referansi} içinde $i_{b,k}^{\star}$ $k$ numaralı bataryanın şarj akımı referansını, $v_{b,k}$ batarya uç gerilimini, $I_{\mathrm{CC},k}$ sabit akım sınırını, $V_{\mathrm{CV},k}$ sabit gerilim eşiğini, $T_{b,k}$ ölçülen batarya sıcaklığını, $T_{\min,k}$ ve $T_{\max,k}$ izin verilen sıcaklık sınırlarını, $K_{v,k}$ gerilim düzenleyici kazancını ve $\operatorname{sat}[\cdot]$ doyum işlemini göstermektedir. Sayısal sınırlar batarya kimyası ve üretici veri sayfasından alınarak atanır.

Coulomb sayımına dayalı ayrık zamanlı şarj durumu güncellemesi
\begin{equation}
  \operatorname{SoC}_k[n+1]=\operatorname{sat}\!\left[
  \operatorname{SoC}_k[n]+\frac{\eta_{c,k}i_{b,k}[n]T_s}{Q_{n,k}},\,0,\,1\right]
  \label{eq:gt-soc-guncelleme}
\end{equation}
olarak yazılmıştır \citep{Plett2015}. Eşitlik~\eqref{eq:gt-soc-guncelleme} içinde $\operatorname{SoC}_k[n]$ $n$ örnekleme anındaki şarj durumunu, $\eta_{c,k}$ Coulomb verimini, $i_{b,k}[n]$ işaretli batarya akımını, $T_s$ BYS örnekleme süresini ve $Q_{n,k}$ güncel kullanılabilir kapasiteyi belirtmektedir; akım işareti şarj yönünde pozitif seçilmiştir. Akım sensörü ofseti ve kapasite yaşlanması nedeniyle yalnız Coulomb sayımı uzun süreli kullanımda sürüklenebilir; başlangıç düzeltmesi ve periyodik kestirim güncellemesi ayrıca tanımlanmalıdır.

\section{Mod Geçişleri ve Güvenlik Kısıtları}
\label{sec:gt-gecis-guvenlik}

RL1 konum değiştirirken akım taşıyan bir bağlantının açılması röle kontaklarında ark ve LCL filtre endüktanslarında aşırı gerilim oluşturabilir; farklı gerilimdeki düğümlerin birleştirilmesi ise sınırlanmamış dengeleme akımına yol açar. Kip geçişi için PWM'nin durdurulması, faz akımlarının sıfıra yaklaşması ve DA bara yarılarının dengede olması birlikte aranmalıdır. Önerilen izin koşulu
\begin{equation}
  \gamma_{k\rightarrow j}=1
  \iff
  \max_{x}|i_{x}|\leq I_{\mathrm{tr}}
  \ \land\
  \bigl|v_{\mathrm{dc}}^{+}-v_{\mathrm{dc}}^{-}\bigr|\leq \Delta V_{\mathrm{tr}}
  \ \land\
  g_{\mathrm{PWM}}=0
  \ \land\
  h_{\mathrm{RL1}}=1
  \label{eq:gt-gecis-izni}
\end{equation}
şeklinde ifade edilmiştir. Eşitlik~\eqref{eq:gt-gecis-izni} içinde $\gamma_{k\rightarrow j}$ $k$ kipinden $j$ kipine geçiş iznini, $i_{x}$ $x\in\{a,b,c\}$ fazının akımını, $I_{\mathrm{tr}}$ geçiş için izin verilen akım eşiğini, $v_{\mathrm{dc}}^{+}$ ve $v_{\mathrm{dc}}^{-}$ bölünmüş DA baranın üst ve alt yarı gerilimlerini, $\Delta V_{\mathrm{tr}}$ izin verilen gerilim farkını, $g_{\mathrm{PWM}}=0$ bütün kapı işaretlerinin durdurulduğunu ve $h_{\mathrm{RL1}}=1$ röle konum geri bildiriminin beklenen konumu doğruladığını gösterir. Sayısal eşikler röle kontak değerleri, yarı iletken güvenli çalışma alanı, filtre endüktansları ve ölçüm belirsizliğiyle belirlenmelidir.

\begin{figure}[H]
  \centering
  \begin{tikzpicture}[
    scale=0.95,
    transform shape,
    node distance=9mm and 13mm,
    state/.style={draw,rounded corners,align=center,minimum width=27mm,minimum height=10mm,fill=blue!5,font=\small},
    safe/.style={state,fill=green!10},
    fault/.style={state,fill=red!8},
    arr/.style={-{Latex[length=2mm]},thick},
    guard/.style={font=\scriptsize,align=center}
  ]
    \node[safe] (off) {Güvenli kapalı};
    \node[state,right=of off] (pre) {Ön şarj ve\\öz sınama};
    \node[state,right=of pre] (ready) {Hazır};
    \node[state,below left=of ready] (charge) {Şarj\\PFC + CC/CV};
    \node[state,below right=of ready] (dis) {Deşarj\\evirici kipi};
    \node[state,below=of ready] (reconf) {RL1\\kip geçişi};
    \node[fault,below=19mm of reconf] (fault) {Kilitli hata};

    \draw[arr] (off) -- node[above,guard]{başlat} (pre);
    \draw[arr] (pre) -- node[above,guard]{ölçüm ve\\geri bildirim uygun} (ready);
    \draw[arr] (ready) -- node[left,guard]{şarj izni} (charge);
    \draw[arr] (ready) -- node[right,guard]{şebeke hazır} (dis);
    \draw[arr] (charge) -- node[above,guard]{akım sıfıra yakın} (reconf);
    \draw[arr] (dis) -- node[above,guard]{akım sıfıra yakın} (reconf);
    \draw[arr] (reconf) -- node[right,guard]{röle\\doğrulandı} (ready);
    \draw[arr] (pre) |- (fault);
    \draw[arr] (charge) -- (fault);
    \draw[arr] (dis) -- (fault);
    \draw[arr] (reconf) -- (fault);
    \draw[arr] (fault) -| node[below,guard]{enerji kesme ve\\yetkili sıfırlama} (off);
  \end{tikzpicture}
  \caption{GT-tipi güvenli durum geçişleri.}
  \label{fig:gt-durum-makinesi}
\end{figure}

Şekil~\ref{fig:gt-durum-makinesi}, çalışma kipini doğrudan değiştirmek yerine önce enerjiyi güvenli seviyeye indiren ve RL1 geçişini ayrı bir durum olarak yürüten sıralamayı göstermektedir. Hızlı yarı iletken koruması donanım düzeyinde çevrimden çevrime çalışmalı; üst seviye durum makinesi röle, kontaktör, sıcaklık, izolasyon ve batarya uygunluk kararlarını yönetmelidir. Yazılımın hata algılama hızı, röle ve kontaktör açma süresi ve DA bara boşalma süresi aynı güvenlik zaman çizelgesinde değerlendirilmelidir.

\begin{table}[H]
  \centering
  \caption{GT-tipi mimari için asgari koruma işlevleri ve doğrulama kanıtları.}
  \label{tab:gt-koruma-dogrulama}
  \small
  \begin{tabularx}{\textwidth}{@{}p{0.22\textwidth}p{0.28\textwidth}X@{}}
    \toprule
    Tehlike / işlev & Gereken koruma & Kabul kanıtı \\
    \midrule
    DA bara aşırı gerilimi & Donanımsal kesme, fren/yük yolu ve yazılımsal sınır & Eşik toleransı, kesme süresi ve enerji boşalma kaydı \\
    Faz/batarya aşırı akımı & Hızlı yarı iletken koruması, sigorta/şalter koordinasyonu & Kısa devre veya güvenli düşük gerilim enjeksiyon testi \\
    Batarya aşırı/düşük gerilimi & Batarya bazlı sınır ve CC--CV sonlandırma & Her kanalın kalibre edilmiş gerilim kaydı ve hata enjeksiyonu \\
    Aşırı/düşük sıcaklık & Şarj/deşarj gücü azaltma ve kilitleme & Sensör kopuk/kısa devre deneyi ve termal eşik kaydı \\
    İzolasyon hatası & Yalıtım izleme, kontaktör açma ve tekrar başlatma kilidi & İzolasyon direnci eşiği ve açma süresi \\
    Hatalı kip geçişi & PWM durdurma, akım sıfırı denetimi, röle geri bildirim kilidi & Kapı/röle zaman diyagramı ve ark/aşırı akım gözlemi \\
    Haberleşme veya işlemci hatası & Güvenli varsayılan durum, bağımsız donanım kesmesi, watchdog & Haberleşme kesme ve watchdog hata enjeksiyonu \\
    Tek batarya arızası & İlgili buck kolunun durdurulması ve şarj gücünün düşürülmesi & Arıza senaryosu boyunca kol akımı, gerilim ve sıcaklık kaydı \\
    \bottomrule
  \end{tabularx}
\end{table}

IEC standartları tek başına belirli bir devrenin güvenli olduğunu kanıtlamaz. Tablo~\ref{tab:gt-koruma-dogrulama}, ilgili standartların kapsamından türetilen tez düzeyindeki asgari kanıt listesidir; nihai uygunluk değerlendirmesi yetkili test, ayrıntılı risk analizi ve ürüne uygulanabilir bütün standartların seçilmesini gerektirir \citep{IEC62619,IEC62477,IEC62933_5_2}.

\section{Klasik Mimari ile Karşılaştırma}
\label{sec:gt-karsilastirma}

Geleneksel EDS mimarisinde şebeke dönüştürücüsü, merkezi DA bara, batarya dizileri ve dengeleme devreleri çoğunlukla ayrı işlevsel katmanlardır. Ayrı katlar tasarım ve servis sınırlarını açık tutarken ilave DA/DA dönüştürücü, dengeleyici, manyetik eleman ve kablolama gerektirebilir. Yeniden yapılandırılabilir sistemler seçilebilir bağlantılarla arıza yalıtımı ve enerji paylaşımı sağlayabilse de her batarya/modül için ek anahtar ve sürücü kullanımı nedeniyle iletim kaybı ile kontrol karmaşıklığı doğurur \citep{Momayyezan2016,Farakhor2023,Huang2024}.

\begin{table}[H]
  \centering
  \caption{GT-tipi mimarinin geleneksel ve literatürdeki yeniden yapılandırılabilir EDS yaklaşımlarıyla nitel karşılaştırması.}
  \label{tab:gt-mimari-karsilastirma}
  \footnotesize
  \begin{tabularx}{\textwidth}{@{}p{0.19\textwidth}p{0.22\textwidth}p{0.25\textwidth}X@{}}
    \toprule
    Ölçüt & Geleneksel sabit dizi + ayrı dönüştürücüler & Yeniden yapılandırılabilir literatür örnekleri & GT-tipi mimari \\
    \midrule
    Şebeke arayüzü & Ayrı iki/üç seviyeli evirici veya doğrultucu & Çalışmaya göre ayrı veya bütünleşik dönüştürücü & Aynı faz kollarında T-tipi evirici (deşarj) ve nötr IGBT'li Vienna PFC (şarj) \\
    Batarya şarjı & Dizi veya modül düzeyinde ortak DA/DA & Modül/hücre seçimi ve enerji yönlendirme & Batarya başına izole doğrultucu--buck kolu \\
    Dengeleme & Pasif/aktif ayrı devre gerekebilir & Anahtar çizelgelemesiyle eşitleyicisiz veya bütünleşik dengeleme mümkün & Şarjda batarya başına bağımsız akım/gerilim düzenlemesi \\
    Arızalı birimi ayırma & Kontaktör veya servis müdahalesi; topolojiye bağlı & Birçok yapıda modül/hücre bypass olanağı & Şarjda ilgili buck kolu durdurulur \\
    Galvanik yalıtım & Merkezi veya modüler DA/DA katında sağlanabilir & Topolojiye göre var/yok & Çok sekonderli BYS transformatörüyle batarya kollarında \\
    Yarı iletken/kontrol yükü & İşlevsel kat sayısına bağlı & Yeniden yapılandırma anahtarlarıyla artabilir & Ayrı PFC ve şarj katı yerine aynı IGBT'ler RL1 ile yeniden kullanılır; röle, diyot ve buck kolları eklenir \\
    Kanıt düzeyi & Olgun ticari ve akademik örnekler mevcut & Benzetim ve ölçekli prototip sonuçları mevcut \citep{Momayyezan2016,Farakhor2023,Huang2024} & Evirici kipinde düşük güç deneyi; şarj yönünde benzetim kayıtları \\
    \bottomrule
  \end{tabularx}
\end{table}

Tablo~\ref{tab:gt-mimari-karsilastirma}, GT-tipi yapının sunduğu olanakları geleneksel ve yeniden yapılandırılabilir yaklaşımlarla yan yana koymaktadır. Mimari, üç işlevi (şebekeye deşarj, şebekeden şarj ve batarya bazlı şarj denetimi) aynı yarı iletken tabanında toplamakta ve ayrı PFC ile şarj dönüştürücüsü yerine aynı IGBT'leri kullanma olanağı sunmaktadır. Başlıca bedeller ise çok sayıda yalıtılmış kanal, sürücü ve algılayıcı; RL1 rölesi ile doğrultucu/boost diyotlarının kayıpları; kip geçiş sırası ve çok sekonderli transformatörün üretilebilirliğidir.

Nitel karşılaştırmayı sayısal düzeye taşıyacak ölçütler de tanımlanmıştır: yarı iletken sayısı, sürücü sayısı, manyetik eleman kütlesi, toplam iletim/anahtarlama kaybı, birim başına ölçüm kanalı ve ortak mod/izolasyon gereksinimleri. Ölçütler, ayrı PFC ve şarj dönüştürücülü bir eşdeğerle aynı güç (\SI{20}{\kilo\volt\ampere}, \num{20}--\SI{25}{\kilo\watt}) ve gerilim tabanında uygulanır.

\section{Doğrulama Durumu ve Mevcut Kanıt Düzeyi}
\label{sec:gt-dogrulama}

GT-tipi mimarinin kanıtı iki gruptadır. Deşarj yönü, yani evirici kipi, Bölüm~\ref{chap:deney} altındaki düşük güç deneyleri ve Bölüm~\ref{chap:anahtarlama}--\ref{chap:sin3} altındaki S1--S4 benzetimleriyle incelenmiştir. Şarj yönü ise Şekil~\ref{fig:gt-viyana-pfc}--\ref{fig:gt-pfc-modu} üzerindeki topoloji ve Şekil~\ref{fig:gt-simulink-modeli}--\ref{fig:gt-batarya-soc} üzerindeki benzetim kayıtlarıyla gösterilmiştir. İki kanıt grubu, bölümün çıktısını bir \emph{mimari tanım, kanıt kaydı ve doğrulama yöntemi} bütünü olarak ortaya koymaktadır. Tablo~\ref{tab:gt-dogrulama-matrisi}, her doğrulama katmanında gösterilen işlevi ve dayandığı kanıt türünü vermektedir.

\begin{table}[H]
  \centering
  \caption{GT-tipi mimari doğrulama matrisi: gösterilen işlevler ve kanıt türleri.}
  \label{tab:gt-dogrulama-matrisi}
  \small
  \begin{tabularx}{\textwidth}{@{}p{0.22\textwidth}p{0.36\textwidth}X@{}}
    \toprule
    Doğrulama katmanı & Gösterilen işlev/çıktı & Kanıt türü \\
    \midrule
    Bağlantı ve durum denetimi & Her mod için iletim yolu, yasak durumlar, kapı/röle/kontaktör sırası & Tasarım tanımı: topoloji şemaları ve durum makinesi. \\
    Evirici kipi (deşarj) & Düşük güçte verim, akım THD'si, güç faktörü & Deney (S1--S3) ve kalibre benzetim (Bölüm~\ref{chap:deney}). \\
    Evirici kipi, anma akımı & Kapalı çevrim akım izleme ve kararlılık & Benzetim ve analitik model: S4 kapalı çevrim modeli ve kararlılık analizi (Bölüm~\ref{chap:sin3}). \\
    Şarj zinciri & PFC akımı, DA bara, primer/sekonder dalga biçimleri, SoC & Benzetim kaydı (Alt Bölüm~\ref{sec:gt-benzetim}). \\
    Batarya/BYS denetimi & SoC güncellemesi, sıcaklık sınırları, CC--CV algoritması & Tasarım tanımı: genel yöntem ve denklemler. \\
    Güvenlik ve koruma & Aşırı akım/gerilim/sıcaklık, izolasyon, sensör ve haberleşme hatalarına karşı koruma işlevleri & Tasarım tanımı: koruma işlevleri ve kabul testi planı. \\
    \bottomrule
  \end{tabularx}
\end{table}

Şarj ve deşarj yönündeki dönüştürme verimleri ayrı ölçülürse gidiş--dönüş verimi
\begin{equation}
  \eta_{\mathrm{rt}}=\eta_{\mathrm{ch}}\eta_{\mathrm{dis}}
  =\frac{E_{b,\mathrm{in}}}{E_{g,\mathrm{in}}}
   \frac{E_{g,\mathrm{out}}}{E_{b,\mathrm{out}}}
  \label{eq:gt-roundtrip-verim}
\end{equation}
olarak tanımlanabilir. Eşitlik~\eqref{eq:gt-roundtrip-verim} içinde $\eta_{\mathrm{ch}}$ şebekeden bataryaya şarj verimini, $\eta_{\mathrm{dis}}$ bataryadan şebekeye deşarj verimini, $E_{g,\mathrm{in}}$ şebekeden çekilen enerjiyi, $E_{b,\mathrm{in}}$ bataryada depolanan enerjiyi, $E_{b,\mathrm{out}}$ bataryadan çekilen enerjiyi ve $E_{g,\mathrm{out}}$ şebekeye verilen enerjiyi göstermektedir. Yardımcı tüketimlerin sistem sınırına dahil edilip edilmediği her sonuçla birlikte belirtilmelidir.

\section{Bölüm Sonucu}

GT-tipi mimari, \num{20}--\SI{25}{\kilo\watt} tasarım hedefli T-tipi deney düzeneğinin RL1 rölesi ve doğrultucu/boost diyotlarıyla şebekeye deşarj, şebekeden şarj ve batarya bazlı şarj denetimi işlevlerine yeniden görevlendirilmesi olarak tanımlanmıştır. Vienna PFC, çok sekonderli BYS, RL1 tabanlı ana topoloji ve üç çalışma yolu devre düzeyinde açıklanmış; şarj zinciri benzetim kayıtlarıyla nitel olarak gösterilmiştir. Evirici kipi, tezin düşük güç deneyleri ve S1--S4 benzetimleriyle incelenen kiptir.

Çalışma modları; güvenli kapalı, ön şarj, PFC/CC--CV şarj, şebekeye deşarj ve RL1 kip geçişi olarak ayrılmıştır. CC--CV referansı, SoC güncellemesi, güvenli kip geçişi ve gidiş--dönüş verimi için kullanılacak denklemler tanımlanmış; bütün parametreler ilk geçtikleri yerde açıklanmıştır. Güvenli geçiş için PWM'nin durdurulması, sıfıra yakın faz akımı, DA bara dengesi ve röle geri bildirimi birlikte şart koşulmuştur.

GT-tipi mimari, T-tipi deney düzeneğinin yarı iletken tabanını şarj ve deşarj işlevlerinde ortak kullanan, tanımlı çalışma kipleri ve güvenli geçiş kurallarıyla donatılmış bir enerji depolama mimarisi sunmaktadır. Evirici kipindeki deneyler ve şarj zinciri benzetim kayıtları, aynı donanımın iki enerji yönünde de işlevini yerine getirdiğini göstermiştir. Tanımlanan gidiş--dönüş verimi ve mimari karşılaştırma ölçütleri, GT-tipi yapının ayrı dönüştürücülü eşdeğerlerle aynı güç tabanında değerlendirilmesine yöntemsel bir temel oluşturmaktadır.
